% GPT-6 Astra Pro; /root/fixed_surface; 22 November 2026 UTC.
% Effort not exposed. Separate proof, awaiting whole independent review.
\section[HD-F1-S36: coherent modules on ordinary powers]{HD-F1-S36: regular normal traces for coherent modules generically free on an ordinary-power thickening}
\label{sec:f1-norm-fourfold-generically-free-powers}

Retain the actual normal integral norm image
$i:Y\hookrightarrow P$, its ideal $\mathcal I$,
the ordinary powers $Y_k$ of S35, and
$E=i_*\mathcal O_Y$ on the S29 geometric open.
Let $M$ be an actual coherent $\mathcal O_{Y_k}$-module
whose localization at the generic point $\xi$ of $Y$
is free of positive rank $m$:
\[
 M_\xi\simeq\mathcal O_{Y_k,\xi}^{\oplus m},
 \qquad F=i_{k*}M,\qquad r_k=\binom{k+3}{4}.
 \tag{F1S36.1}
\]
There is no local-freeness or purity requirement
elsewhere on $Y_k$. In particular, all actual
lower-dimensional torsion and module extensions
in $M$ are retained.
Put $D=\mathfrak D_P$ and
$W_E=\operatorname{ob}_E(D)$.
The result is the actual characteristic-image
surjection
\[
\begin{split}
 \mathcal T_{k,M}:
 \operatorname{ev}_F\bigl(D\oplus H^2(P,\mathcal O_P)\bigr)
 &\longrightarrow W_E,\\
 \operatorname{ev}_F(\kappa+\eta)
 &\longmapsto -m r_k\operatorname{ob}_E(\kappa).
\end{split}
 \tag{F1S36.2}
\]
For every such $F$ with the complete S14 character
\[
 e^B\operatorname{ch}(F)=\lambda t+Q(\beta),
 \qquad B\in H^{1,1}(P,\mathbf Q),\quad
 0\ne\lambda t\in W_K(P),\quad Q\in\mathbf Q[x],
 \tag{F1S36.3}
\]
its invariant full degree-two trace kernel has
dimension at least $16-r_C\ge4$ for every specified
coherent finite point lift. Thus the stated
coherent ordinary-power class fails the actual
weak semiregularity gate.

\paragraph{1. The generic ideal action and its canonical normal value.}
Suppose $\operatorname{ev}_F(\kappa+\eta)=0$.
Use precisely the Toda/BF sign and deformation
comparison retained in S33--S35. The actual
perfect lift has central fibre the sheaf $F$,
so the S33 cohomology argument gives an actual
flat sheaf $\widetilde M$ on each commutative
local first-order ambient deformation. Its
overlap maps retain the prescribed scalar
twisting. Fix one such ambient trivialization
and write its ring as $R_\epsilon=R[\epsilon]/(\epsilon^2)$.

For $f\in I^k$, the central module is annihilated
by $f$, so
\[
 f\widetilde M\subset\epsilon\widetilde M
       \simeq M,\qquad
 a_f(v)=\epsilon^{-1}f\widetilde v.
 \tag{F1S36.4}
\]
Here $\epsilon^{-1}$ denotes the displayed
isomorphism with the epsilon submodule, not
division in a ring. A change of lift of $v$
adds $\epsilon w$, whose contribution is zero
because $fM=0$. Thus $a_f$ is an actual
$R/I^k$-linear endomorphism of $M$, and
$f\mapsto a_f$ is $R$-linear.

There is a dense open $U_0\subset Y_{\rm reg}$
on which $M$ is locally free over
$\mathcal O_{Y_k}$. This follows by clearing
the finitely many denominators in the generic
isomorphism (F1S36.1) and its inverse.
On $U_0$ form the ordinary matrix trace
\[
 \delta(f)=\frac1m\operatorname{Tr}_{\mathcal O_{Y_k}}(a_f),
 \qquad
 u=\rho_k(\delta):\mathcal I/\mathcal I^2
                         \longrightarrow\mathcal O_{U_0}.
 \tag{F1S36.5}
\]
The map $\rho_k$ is the actual S35 symmetric-power
trace map. Both maps are invariant under change
of the local free basis. In particular $u$ gives
a canonically specified rational conormal
functional on $Y$ for this actual deformation.
The next three paragraphs prove that each of
its values extends regularly; no extension of
the possibly singular nonreduced matrix trace
$\delta$ is assumed.

\paragraph{2. The exact moment calculation over a field.}
Let $L$ be a characteristic-zero field,
$S=L[[x_1,\ldots,x_c]]$, $c\ge2$,
$\mathfrak n=(x_1,\ldots,x_c)$ and
$A=S/\mathfrak n^k$.
Take an actual $L[\epsilon]$-flat finite module
$N_\epsilon$ over $S[\epsilon]$ whose central
module is $A^{\oplus m}$.
It is finite free over $L[\epsilon]$: a lift
of a central $L$-basis generates by
$\epsilon^2=0$, and flatness makes the resulting
kernel have zero central fibre and hence zero.

Choose lifts $\widetilde e_1,\ldots,\widetilde e_m$
of the standard $A$-basis. The actual vectors
\[
           x^\alpha\widetilde e_a,
   \qquad |\alpha|<k,\quad 1\le a\le m,
 \tag{F1S36.6}
\]
are an $L[\epsilon]$-basis because their reductions
are the standard monomial basis.
For $f\in\mathfrak n^k$, (F1S36.4) now gives
\[
 a:\mathfrak n^k\longrightarrow
                         \operatorname{Mat}_m(A).
 \tag{F1S36.7}
\]
Each matrix entry is an actual $S$-linear map.
Apply the complete S35 monomial proof
to each entry, using (F1S35.5)--(F1S35.8). It gives
\[
 a(\mathfrak n^k)\subset
       \operatorname{Mat}_m(\mathfrak n^{k-1}/\mathfrak n^k),
 \qquad a(\mathfrak n^{k+1})=0.
 \tag{F1S36.8}
\]
No splitting of $N_\epsilon$ as a deformed
sum of modules is used here.

Multiplication by $x_i$ on (F1S36.6) is the
literal monomial shift when $|\alpha|<k-1$.
When $|\alpha|=k-1$, it is
$\epsilon a_{x_i x^\alpha}(e_a)$.
The first kind has no diagonal coefficient.
The trace of the second kind is exactly
\[
\begin{split}
 \operatorname{Tr}_{L[\epsilon]}(\mu_{x_i})
 &=\epsilon\sum_{|\alpha|=k-1}\sum_{a=1}^m
       [x^\alpha](a_{x_i x^\alpha})_{aa}\\
 &=\epsilon m r_k\,\rho_k(\delta)(x_i),
 \qquad
 \delta=\frac1m\operatorname{Tr}_A(a),
 \qquad r_k=\binom{c+k-1}{c}.
\end{split}
 \tag{F1S36.9}
\]
The second equality is the definition of
$\rho_k$: its unnormalized trace is the sum
of the displayed diagonal coefficients on
$\operatorname{Sym}^{k-1}(\mathfrak n/\mathfrak n^2)$.

For a monomial $x^\gamma$ of degree $d\ge2$,
its action on (F1S36.6) has no diagonal term.
If $|\alpha|+d<k$, the degree increases by $d$.
If $|\alpha|+d=k$, its epsilon output has degree
$k-1$ by (F1S36.8), strictly larger than
$|\alpha|=k-d$.
If $|\alpha|+d>k$, its action is zero by the
second assertion of (F1S36.8).
Therefore
\[
 \operatorname{Tr}_{L[\epsilon]}(\mu_f)=0
       \quad(f\in\mathfrak n^2),
 \qquad
 \operatorname{Tr}_{L[\epsilon]}(\mu_f)
       =\epsilon m r_k\,u(f)
       \quad(f\in\mathfrak n).
 \tag{F1S36.10}
\]
These identities apply to power series as well:
the actual action of $\mathfrak n^{2k}$ on
$N_\epsilon$ is already zero, since one
$\mathfrak n^k$ factor lands in epsilon and
the second kills its central module.
Hence only finitely many monomials act.
For $k=1$ the same calculation has the single
monomial of degree zero, and products of two
normal variables act through $\epsilon^2=0$.

\paragraph{3. An actual finite free resolution over the support coefficients.}
Let $y$ be any closed point of $Y_{\rm reg}$
in the chosen ambient chart, including a point
where $M$ is not locally free over $Y_k$.
Use the compatible completed-coordinate maps
of (F1S35.9)--(F1S35.10):
\[
 \widehat R=B[[x_1,x_2,x_3,x_4]],\quad
 B=\mathbf C[[t_1,t_2,t_3,t_4]],\quad
 I\widehat R=(x_1,x_2,x_3,x_4)=\mathfrak n,
 \quad \widehat{\mathcal O}_{Y,y}=B.
 \tag{F1S36.11}
\]
Complete the actual local lifted module by the
flat map $R[\epsilon]\to\widehat R[\epsilon]$,
and denote it by $N_\epsilon$.
Flatness preserves its exact epsilon sequence,
so $N_\epsilon$ is $\mathbf C[\epsilon]$-flat.
Its central module $N$ is finite over
$\widehat R/\mathfrak n^k$, and hence finite
over $B$.

The module $N_\epsilon$ is finite over
$B_\epsilon=B[\epsilon]/(\epsilon^2)$.
Indeed lift finitely many $B$-generators of
$N$ and take their $B_\epsilon$-span.
The quotient $Q$ satisfies $Q=\epsilon Q$,
whence $Q=\epsilon^2Q=0$.
This proves finiteness without first assuming it.

The regular local ring $B$ has dimension four.
Thus $N$ has an actual finite free $B$-resolution
of length at most four, by Stacks 00O7.
That resolution lifts to a finite free
$B_\epsilon$-resolution of $N_\epsilon$ as follows.
Lift the first free module's generators to obtain
an actual surjection $F_{0,\epsilon}\to N_\epsilon$.
Its kernel is finite and epsilon-flat.
Reducing the sequence modulo epsilon is exact:
the periodic free resolution of $B$ over
$B_\epsilon$, with all differentials epsilon,
computes its positive Tor groups as
$\ker\epsilon/\epsilon N_\epsilon=0$.
Repeat with the actual central syzygies.
At the last stage the central syzygy is free
over $B$. Lifting its basis gives a surjection
from a finite free $B_\epsilon$-module.
Its kernel has zero reduction, again by
epsilon-flatness, and is zero by
$Q=\epsilon Q\Rightarrow Q=0$.
This constructs the actual bounded free resolution
\[
      P_\epsilon^\bullet\longrightarrow N_\epsilon.
 \tag{F1S36.12}
\]
The assertion of epsilon-flatness of the
successive kernels follows from the exact
sequence of Tor over $\mathbf C[\epsilon]$,
since the free middle module and the quotient
are both flat. No ideal-filtration quotient of
$M$ or $\widetilde M$ was required to be flat.

For each actual multiplication endomorphism
$\mu_f$ of $N_\epsilon$, lift it to a chain map
of (F1S36.12), and take
\[
       \operatorname{tr}_{B_\epsilon}(\mu_f)
        =\sum_j(-1)^j\operatorname{Tr}(\mu_f^j)
                                      \in B_\epsilon.
 \tag{F1S36.13}
\]
This is the categorical trace of the actual
perfect $B_\epsilon$-module $N_\epsilon$.
The matrix expression is unchanged by a chain
homotopy: the two sums from $dh+hd$ cancel
by cyclicity of rectangular matrix traces.
Comparison maps between free resolutions,
their homotopies to the identity, and the same
cyclicity prove independence of the resolution.
Equivalently it is the actual dual, coevaluation,
Koszul symmetry and evaluation composite of S32.
The finite matrix expression commutes with
base change. In particular it commutes with
$B_\epsilon\to L_\epsilon$, where
$L=\operatorname{Frac}(B)$.

After this localization, the central module is
$\bigl(L[[x_1,\ldots,x_4]]/\mathfrak n^k\bigr)^m$
in the sense of its finite monomial quotient.
Indeed the dense free open $U_0$ pulls back
to that generic coefficient point: every
nonzero support denominator stays nonzero
in the domain $B$ and becomes invertible in $L$.
The lifted module is finite free over
$L_\epsilon$, by its epsilon-flatness and
its finite central dimension. Thus the trace
of the localized resolution is precisely the
ordinary trace computed in (F1S36.9)--(F1S36.10).
The formal-power-series notation introduces
no completion/localization interchange:
the action factors through the finite quotient
by $\mathfrak n^{2k}$ throughout.

\paragraph{4. Descent of the regular normal value.}
For $f\in I_y$, its image in the completed
ring belongs to $\mathfrak n$.
Equation (F1S36.10) and the injectivity
$B_\epsilon\hookrightarrow L_\epsilon$ show that
the central coefficient of (F1S36.13) is zero.
Consequently there is a unique $b_f\in B$ with
\[
       \operatorname{tr}_{B_\epsilon}(\mu_f)
                     =\epsilon m r_k b_f.
 \tag{F1S36.14}
\]
Under the actual inclusion
$\mathbf C(Y)\hookrightarrow\operatorname{Frac}(B)$,
equation (F1S36.9) identifies $b_f$ with the
rational value $u(f)$ of (F1S36.5).
This uses the literal base change of the S35
conormal multiplication and symmetric trace;
both commute with the displayed flat maps.
It follows in particular that the value is
independent of the chosen completed parameters
and of the coefficient ring used to verify it.

The map $A_0=\mathcal O_{Y,y}\to B=\widehat A_0$
is faithfully flat. Write $u(f)=a/b$ with
$a,b\in A_0$, $b\ne0$.
Membership of its image in $B$ means
$a\in bB$. Faithful flatness gives
$bB\cap A_0=bA_0$; explicitly, it detects
the vanishing of the class of $a$ in
$A_0/bA_0$ after tensoring with $B$.
Hence $u(f)\in A_0$.
Thus every actual local ideal generator has a
regular conormal value at every closed smooth
point, even where $M$ has arbitrary defects.

Finite generation of the ideal makes these
values regular on neighborhoods of those
points. Such neighborhoods cover the smooth
locus, since a nonempty closed complement in
a finite-type complex scheme has a closed point.
On the dense free open the values are linear,
annihilate $I^2$, and satisfy all original
conormal relations. The same equations hold
on the smooth locus by integrality.
Finally every height-one point of the normal
$Y$ lies in that locus. The normal-domain
intersection of Stacks 0AVB extends each of
the finitely many generator values through
the remaining codimension at least two set.
Every relation again vanishes by density.
This constructs the actual regular map
\[
                    u:\mathcal I/\mathcal I^2
                                       \longrightarrow\mathcal O_Y
 \tag{F1S36.15}
\]
on the ambient chart under consideration.
Only the reduced support is used for this
normal extension. There is no assertion that
the generic trace with values in
$\mathcal O_{Y_k,\xi}$ extends to $Y_k$.

\paragraph{5. Scalar-twisted overlap maps and the embedded deformation.}
On an overlap of two ambient trivializations,
restrict first to the dense open where $M$ is
free over $Y_k$.
The actual ideal-action calculation of S35 gives
\[
        \delta_i-\delta_j=\partial_k(\kappa_{ij}).
 \tag{F1S36.16}
\]
Indeed change of module trivialization conjugates
the endomorphism, while the ambient transition
adds multiplication by the derivative of the
ideal element. The normalized matrix trace
of that scalar is the derivative itself.
The scalar triple products specified by $\eta$
act trivially by conjugation, so this comparison
retains the given coherent twisted gluing.
Apply the exact S35 identity
$\rho_k\partial_k=\partial_1$ and then use the
regular extensions just constructed. Equality
on the dense open gives equality everywhere:
\[
                  u_i-u_j=\partial_1(\kappa_{ij}).
 \tag{F1S36.17}
\]

On each chart use the explicit ideal
\[
 I_{\epsilon,i}=\{f+\epsilon h: f\in I,
                         \ h\bmod I=-u_i(f)\}.
 \tag{F1S36.18}
\]
Linearity and annihilation of $I^2$ make it an
ideal. Its quotient has
$\ker\epsilon=\operatorname{im}\epsilon$:
an element killed by epsilon has constant part
in $I$, and subtracting its prescribed lift in
(F1S36.18) leaves an epsilon class; conversely
epsilon squares to zero.
The quotient is therefore flat over the dual
numbers. Equation (F1S36.17) gives its actual
overlap compatibility. We obtain an embedded
flat deformation of the original reduced
$Y$ in the given ambient deformation.
The exact S24 ideal-obstruction comparison yields
\[
 \operatorname{ev}_F(\kappa+\eta)=0
                 \quad\Longrightarrow\quad
                         \operatorname{ob}_E(\kappa)=0.
 \tag{F1S36.19}
\]
This is a consequence of the constructed ideals,
not of a proposed extension of the support
filtration of $\widetilde M$.

\paragraph{6. The actual characteristic image and full trace.}
Subtracting two Hochschild representatives with
the same evaluation and using (F1S36.19) proves
that (F1S36.2) is well defined. It is linear.
Its image is $W_E$, since this is the image of
the actual ambient obstruction map and
$m r_k\ne0$.
The minus sign is the corrected first-jet/HKR
dictionary of S31, retained in S35.

At the generic point the module has length
$m r_k$ over the regular ambient local ring
of dimension four. The actual localized
leading-cycle map gives
\[
 \operatorname{ch}_j(F)=0\quad(j<4),\qquad
 \operatorname{ch}_4(F)
             =m r_k[Y]=m r_k\left(w+\frac{\beta^4}{72}\right)
                   \quad\hbox{in rational cohomology}.
 \tag{F1S36.20}
\]
All smaller-dimensional torsion and defects
remain in the higher components; none is
discarded from (F1S36.3).
Comparing the complete S14 equality in degree
four, then cupping with $\beta$, fixes its
polynomial coefficient and gives
\[
              q_4=\frac{m r_k}{72},\qquad
                         \lambda t=m r_k w.
 \tag{F1S36.21}
\]
Here $\beta w=0$ and $\beta^5\ne0$ are the
retained actual norm-image identities.
The actual sheared characteristic image
\[
 \operatorname{ev}_F
       \{\kappa+\kappa\mathbin{\lrcorner}B:\kappa\in D\}
 \tag{F1S36.22}
\]
has zero full trace by the complete BF
exponential derivative used in S31--S35.
Its image under (F1S36.2) is $W_E$.
Thus its dimension, and hence the full trace
kernel dimension, is at least
$\dim W_E=16-r_C\ge4$.
Every specified coherent finite point lift
fixes this image by naturality of the actual
Hochschild action; translations and
$\operatorname{Pic}^0$ tensor functors act
trivially on the degree-two source.
The exact R05 equivariant trace comparison
therefore gives the stated invariant kernel.
Since $F$ is an actual sheaf, its negative
self-Ext groups vanish.

\paragraph{7. Concrete additional modules and the trace boundary.}
There are actual modules of the present class
which are not vector bundles on $Y_k$.
Take a closed point $p\in Y_{\rm reg}$ and the
literal kernel
\[
        0\longrightarrow M_p\longrightarrow
              \mathcal O_{Y_k}\longrightarrow\mathbf C_p
                                      \longrightarrow0.
 \tag{F1S36.23}
\]
Its generic module is $\mathcal O_{Y_k,\xi}$.
At $p$ it is the maximal ideal of the local
ring of $Y_k$. For $k\ge2$ that ring has
cotangent dimension eight, since
$I^k\subset\mathfrak m_P^2$; its maximal
ideal consequently needs eight generators.
A locally free module of generic rank one
would need one. Thus $M_p$ is not locally
free there. For $k=1$ the same argument uses
the four-dimensional smooth support.
For $k\ge2$, the generic action of an ideal
element with nonzero conormal class on $M_p$
is nonzero, as it is on $\mathcal O_{Y_k}$.
The explicit ideal-action comparison of S35
therefore also compares these modules with
the reduced-support scope of S34.

More generally, for an actual vector bundle
$V$ of rank $m$ on $Y_k$, any surjection of
its fibre at $p$ onto $\mathbf C$ gives a
kernel $M_{V,p}$ of the present class.
Locally that kernel is
$\mathfrak m_{Y_k,p}\oplus\mathcal O_{Y_k,p}^{m-1}$,
so the same generator count proves its
failure of local freeness.
The exact coherent sequence gives the actual
$K$-class difference $[i_{k*}V]-[i_{k*}M_{V,p}]
=[\mathbf C_p]$.
After the cycle-class map to rational cohomology,
\[
 \operatorname{ch}(i_{k*}M_{V,p})
     =\operatorname{ch}(i_{k*}V)-[p],\qquad
 [p]=\frac{\beta^8}{\int_P\beta^8},\qquad e^B[p]=[p].
 \tag{F1S36.24}
\]
Hence the full S14 form for these two actual
objects is equivalent, with only the displayed
top polynomial coefficient changed.
No rational equivalence of different points
in the Chow group is asserted.

The use of the regular coefficient ring in
paragraphs 3--4 is essential to the proved
trace construction. For an exact local
illustration set
\[
 B_0=\mathbf C[[t]],\quad
 A_0=B_0[x_1,x_2,x_3,x_4]/(x_1,x_2,x_3,x_4)^2,
 \quad M_0=(t,x_1,x_2,x_3,x_4)\subset A_0.
 \tag{F1S36.25}
\]
The module is generically free of rank one.
Multiplication by $x_1/t$ in its generic
localization preserves $M_0$: it sends $t$
to $x_1$ and every $x_j$ to zero.
It therefore defines an actual $A_0$-linear
endomorphism. Its generic $A_0$-trace is
$x_1/t$, which is not an element of $A_0$.
This verifies that arbitrary nonreduced
generic matrix traces need not extend, even
in the present generic module class.
The proof instead extends the specific
normal value (F1S36.5), by its exact
regular-trace identity (F1S36.14).
No assertion is made here for generic
modules without the freeness (F1S36.1),
or for arbitrary multi-cohomology objects
over the nonreduced $Y_k$.

\paragraph{Primary-source record.}
The monomial calculation, symmetric trace,
completed coordinate maps and extension
across the normal support are the exact
S35 constructions. The field-valued moment
calculation and the actual epsilon-flat
syzygy lifting above are proved directly.
The finite free resolution over the regular
coefficient ring uses Stacks 00O7, with its
complete supplied proof and the complete
used 00NT, 00NG, 00NS, 00NE and 00N6
proofs freshly read.
The actual perfect trace and its signs use
the S32 composites; complete 0FNK, 08DQ
and 0FVB proofs were freshly reread.
The coefficient-ring trace is only applied
to the actual perfect module (F1S36.12).
The Toda lift, flat-sheaf extraction,
full-character derivative and finite-action
comparison retain their precise S33--S35
source scopes.
